\section{Spectral fitting and photon path-length derivation}
\label{app:spectral_fit_derivation}

%\subsection{}     %% Appendix A1, A2, etc.


The high-resolution vertical optical depth is computed from layer extinction as
\begin{equation}
\tau_v(\lambda_q)
=
\sum_z
\alpha(\lambda_q,z)\Delta z
+
\tau_{\mathrm{Ray}}(\lambda_q),
\label{eq:vertical_tau_hr}
\end{equation}
where $\alpha$ includes the target absorber and H$_2$O absorption on the ABSCO
grid, and $\tau_{\mathrm{Ray}}$ is the Rayleigh optical depth.


The high-resolution slant optical depth is
\begin{equation}
\mathrm{SOD}(\lambda_q)
=
\left(
\frac{1}{\cos\theta} + \frac{1}{\cos\phi}
\right)
\tau_v(\lambda_q).
\label{eq:hr_sod}
\end{equation}


Because the instrument line shape operates on transmission rather than optical
depth, the channel value used in the fit is
\begin{equation}
\overline{\mathrm{SOD}}_c
=
-\ln
\left[
\frac{
\sum_q R_{cq}\exp\{-\mathrm{SOD}(\lambda_q)\}
}{
\sum_q R_{cq}
}
\right].
\label{eq:discrete_ils_sod}
\end{equation}


The high-resolution solar spectrum is adjusted for Earth--Sun distance and the
solar Doppler stretch before convolution with the instrument line shape:
\begin{equation}
F_{\odot}^{\ast}(\lambda_q)
=
\frac{
F_{\odot}(\lambda_q)
}{
d_{\mathrm{AU}}^2
}
\frac{1}{1+\beta},
\qquad
\lambda_q^{\ast}
=
\lambda_q(1+\beta),
\label{eq:solar_distance_doppler}
\end{equation}
where $d_{\mathrm{AU}}$ is the Earth--Sun distance in astronomical units and
$\beta=v_{\mathrm{solar}}/c$.


The channel solar term is then computed as the ILS-weighted mean,
\begin{equation}
\mathcal{F}_{0,c}^{\mathrm{ILS}}
=
\frac{
\sum_q R_{cq}F_{\odot}^{\ast}(\lambda_q^{\ast})
}{
\sum_q R_{cq}
}.
\label{eq:solar_ils_convolution}
\end{equation}


\begin{equation}
\mathcal{F}_{0,c}
=
m_1
\mathcal{F}_{0,c}^{\mathrm{ILS}}
\frac{1}{\cos\theta},
\qquad
m_1 = 0.5,
\label{eq:effective_solar_normalization}
\end{equation}


\begin{equation}
T_c
=
\frac{\pi I_c}{\mathcal{F}_{0,c}}.
\label{eq:code_transmittance}
\end{equation}






For a photon traveling along the path coordinate $s$ over the total path length $L$, the Beer--Lambert transmittance is
\begin{equation}
T_{\lambda}(L)
=
\exp\left[
-
\int_0^L
\sigma_{\lambda}(s)n(s)
\,\mathrm{d}s
\right]
=
\exp\left[
-
\int_0^L
\alpha_{\lambda}(s)
\,\mathrm{d}s
\right],
\label{eq:beer_lambert_path}
\end{equation}



where $\alpha_{\lambda}(s)=\sigma_{\lambda}(s)n(s)$ is the absorption
coefficient. Averaging over photons with different path lengths gives
\begin{equation}
\bar{T}_{\lambda}
=
\frac{1}{N}
\sum_{i=1}^{N}
\exp\left[
-
\int_0^{L_i}
\alpha_{\lambda,i}(s)
\,\mathrm{d}s
\right].
\label{eq:ensemble_transmittance}
\end{equation}


Under the mean-absorption approximation,
$\alpha_{\lambda,i}(s)\approx\bar{\alpha}_{\lambda}$, so that
\begin{equation}
\bar{T}_{\lambda}
\approx
\frac{1}{N}
\sum_{i=1}^{N}
\exp(-\bar{\alpha}_{\lambda}L_i)
=
\int_0^\infty
p_L(L)
\exp(-\bar{\alpha}_{\lambda}L)
\,\mathrm{d}L .
\label{eq:path_length_laplace_dimensional}
\end{equation}

\begin{equation}
\overline{\mathrm{SOD}}_j
=
-\ln
\left[
\frac{
\sum_q R_j(\lambda_q)
\exp\{-\mathrm{SOD}(\lambda_q)\}
}{
\sum_q R_j(\lambda_q)
}
\right],
\end{equation}

where $j$ indexes the discrete wavelength samples in an OCO-2 band and
$q$ indexes the high-resolution wavelength grid used for the ILS convolution.


Let $\bar{L}$ be the direct geometric slant path and define the dimensionless
path-length enhancement
\begin{equation}
l' = \frac{L}{\bar{L}},
\qquad
\overline{\mathrm{SOD_{c, \lambda}}}
=
\bar{\alpha}_{\lambda}\bar{L}.
\label{eq:path_enhancement_tau}
\end{equation}



The ensemble transmittance can then be written as the Laplace transform of
the path-enhancement distribution:
\begin{equation}
\bar{T}(\overline{\mathrm{SOD}}_c)
=
\int_0^\infty
p(l')
\exp(-\overline{\mathrm{SOD}}_c \, l')
\,\mathrm{d}l' .
\label{eq:path_enhancement_laplace}
\end{equation}


For a Lambertian surface with reflectance $\gamma$,
\begin{equation}
\pi I_c(\lambda)
=
\mathcal{F}_{0,c}(\lambda)\cos\theta\,
\gamma\,
\bar{T}(\overline{\mathrm{SOD}}_{c \, \lambda}),
\label{eq:radiance_lambertian}
\end{equation}



and therefore
\begin{equation}
T(\overline{\mathrm{SOD}}_c)
\equiv
\frac{\pi I_c}{\mathcal{F}_{0,c}\cos\theta}
=
\gamma
\int_0^\infty
p(l')
\exp(-\overline{\mathrm{SOD}}_c \, l')
\,\mathrm{d}l' .
\label{eq:measured_laplace_appendix}
\end{equation}



If $p(l')$ is approximated by a gamma distribution with mean
$\mu_{l'}=\langle l'\rangle$ and shape parameter
$\kappa=\mu_{l'}^2/\operatorname{var}(l')$, then
\begin{equation}
p(l')
=
\frac{(\kappa/\mu_{l'})^{\kappa}}{\Gamma(\kappa)}
l'^{\kappa-1}
\exp\left(-\frac{\kappa}{\mu_{l'}}l'\right),
\qquad
l' \ge 0 .
\label{eq:gamma_path_distribution}
\end{equation}



Its Laplace transform is
\begin{equation}
\int_0^\infty
p(l')\exp(-\overline{\mathrm{SOD}}_c \, l')
\,\mathrm{d}l'
=
\left(
1+\frac{\mu_{l'}}{\kappa}\overline{\mathrm{SOD}}_c
\right)^{-\kappa},
\label{eq:gamma_laplace_transform}
\end{equation}



so the measured transmittance is
\begin{equation}
T(\overline{\mathrm{SOD}}_c)
=
\gamma
\left(
1+\frac{\mu_{l'}}{\kappa}\overline{\mathrm{SOD}}_c
\right)^{-\kappa}.
\label{eq:gamma_transmittance}
\end{equation}



Taking the logarithm gives
\begin{equation}
\ln T(\overline{\mathrm{SOD}}_c)
=
\ln\gamma
-
\kappa
\ln\left(
1+\frac{\mu_{l'}}{\kappa}\overline{\mathrm{SOD}}_c
\right).
\label{eq:log_gamma_transmittance}
\end{equation}


Using the Taylor expansion of the logarithm,
\begin{equation}
\ln T(\overline{\mathrm{SOD}}_c)
=
\ln\gamma
-
\mu_{l'}\overline{\mathrm{SOD}}_c
+
\frac{\mu_{l'}^2}{2\kappa}\overline{\mathrm{SOD}}_c^2
-
\frac{\mu_{l'}^3}{3\kappa^2}\overline{\mathrm{SOD}}_c^3
+
\cdots .
\label{eq:log_gamma_taylor}
\end{equation}



The fitted polynomial is written as
\begin{equation}
\ln T(\overline{\mathrm{SOD}}_c)
=
b
+
\sum_{n=1}^{N}
(-1)^n
\frac{k_n}{n}
\overline{\mathrm{SOD}}_c^n,
\qquad
b=\ln\gamma .
\label{eq:log_polynomial_fit_appendix}
\end{equation}



For the gamma model, matching coefficients gives
\begin{equation}
k_n
=
\frac{\mu_{l'}^{n}}{\kappa^{n-1}},
\qquad
n\ge 1 .
\label{eq:gamma_coefficient_identity}
\end{equation}




Thus
\begin{equation}
k_1=\mu_{l'},
\qquad
k_2=\frac{\mu_{l'}^{2}}{\kappa}
=
\operatorname{var}(l'),
\qquad
\kappa=\frac{k_1^2}{k_2}.
\label{eq:k1_k2_kappa_identity}
\end{equation}